% GENERATED FILE. Edit canonical lessons, then run npm run generate:books.
% canonical-source-hash: fnv1a64:377c4d6df294122d
% canonical-lessons: MR-C509-padonnati, MR-C509-pad, MR-C509-varistha, MR-C509-kanistha, MR-C509-jhatne

\chapter{Moving Up}
\label{ch:mr-moving-up}
\hlchaptermodality{509}

\begin{chapteropening}
\textbf{By the end of this chapter:} \emph{I can talk about posts, seniors and juniors and moving up at work.}

Complete the last lesson of chapter 509: I can talk about posts, seniors and juniors and moving up at work.
\end{chapteropening}

\section[\texorpdfstring{\mr{पदोन्नती} (padonnatī)}{पदोन्नती (padonnatī)}]{\mr{पदोन्नती} (padonnatī) — a promotion}

\label{lesson:MR-C509-padonnati}



\begin{quote}
Before the new one: say the Marathi for a pay rise, then the Marathi for to earn.
\end{quote}

\subsection*{You'll want to know: \mr{पदोन्नती}}
\textbf{\mr{पदोन्नती}} — \emph{padonnatī} — \textquotedblleft{}a promotion\textquotedblright{}.

\begin{quote}
\mr{मीराला} \mr{या} \mr{वर्षी} \mr{पदोन्नती} \mr{मिळाली}.
\end{quote}

\begin{quote}
\emph{mīrālā yā varṣī padonnatī miḷālī.}
\end{quote}

\begin{quote}
Meera got a promotion this year.
\end{quote}

Feminine.

\begin{cousinweb}
\textbf{\mr{पदोन्नती}} (\emph{padonnatī}) is a tatsama compound of Sanskrit \emph{pada}, \textquotedblleft{}step, position\textquotedblright{}, and \emph{unnati}, \textquotedblleft{}rising\textquotedblright{} (Monier-Williams).
\end{cousinweb}

\begin{grammarlens}[title={What you've built}]
One more word for describing what you lived and felt.
\end{grammarlens}

\subsection*{Your turn}
\begin{itemize}
\raggedright
  \item \emph{Say it:} \emph{padonnatī}
  \item \emph{Say it:} \emph{padonnatī}, once more
  \item \emph{Say it:} \emph{mīrālā yā varṣī padonnatī miḷālī.} — Meera got a promotion this year.
  \item \emph{Recall:} say the Marathi for a pay rise, then the Marathi for to earn, then say \emph{padonnatī} again
\end{itemize}

\begin{tcolorbox}[breakable,colback=teal!4,colframe=teal!35!black,title={Before you move on}]
What does \mr{पदोन्नती} mean? (\textquotedblleft{}A promotion\textquotedblright{}.) Say it once more.
\end{tcolorbox}
\section[\texorpdfstring{\mr{पद} (pad)}{पद (pad)}]{\mr{पद} (pad) — a position; a post}

\label{lesson:MR-C509-pad}



\begin{quote}
Before the new one: say the Marathi for to earn, then the Marathi for a promotion.
\end{quote}

\subsection*{You'll want to know: \mr{पद}}
\textbf{\mr{पद}} — \emph{pad} — \textquotedblleft{}a position; a post\textquotedblright{}.

\begin{quote}
\mr{तिला} \mr{व्यवस्थापकाचं} \mr{पद} \mr{मिळालं}.
\end{quote}

\begin{quote}
\emph{tilā vyavasthāpakācaṁ pad miḷālaṁ.}
\end{quote}

\begin{quote}
She got the manager's post.
\end{quote}

Neuter.

\begin{cousinweb}
\textbf{\mr{पद}} (\emph{pad}) is a tatsama, taken from Sanskrit \emph{pada}, \textquotedblleft{}step, place, position\textquotedblright{} (Monier-Williams).
\end{cousinweb}

\begin{grammarlens}[title={What you've built}]
One more word for describing what you lived and felt.
\end{grammarlens}

\subsection*{Your turn}
\begin{itemize}
\raggedright
  \item \emph{Say it:} \emph{pad}
  \item \emph{Say it:} \emph{pad}, once more
  \item \emph{Say it:} \emph{tilā vyavasthāpakācaṁ pad miḷālaṁ.} — She got the manager's post.
  \item \emph{Recall:} say the Marathi for to earn, then the Marathi for a promotion, then say \emph{pad} again
\end{itemize}

\begin{tcolorbox}[breakable,colback=teal!4,colframe=teal!35!black,title={Before you move on}]
What does \mr{पद} mean? (\textquotedblleft{}A position; a post\textquotedblright{}.) Say it once more.
\end{tcolorbox}
\section[\texorpdfstring{\mr{वरिष्ठ} (variṣṭha)}{वरिष्ठ (variṣṭha)}]{\mr{वरिष्ठ} (variṣṭha) — senior}

\label{lesson:MR-C509-varistha}



\begin{quote}
Before the new one: say the Marathi for a promotion, then the Marathi for a position.
\end{quote}

\subsection*{You'll want to know: \mr{वरिष्ठ}}
\textbf{\mr{वरिष्ठ}} — \emph{variṣṭha} — \textquotedblleft{}senior\textquotedblright{}.

\begin{quote}
\mr{वरिष्ठ} \mr{अधिकारी} \mr{आज} \mr{बैठकीला} \mr{आले}.
\end{quote}

\begin{quote}
\emph{variṣṭha adhikārī āj baiṭhkīlā āle.}
\end{quote}

\begin{quote}
The senior officers came to the meeting today.
\end{quote}

As a noun in the plural: \mr{वरिष्ठ}, one's seniors.

\begin{cousinweb}
\textbf{\mr{वरिष्ठ}} (\emph{variṣṭha}) is a tatsama, taken from Sanskrit \emph{variṣṭha}, \textquotedblleft{}best, foremost\textquotedblright{}, a superlative (Monier-Williams).
\end{cousinweb}

\begin{grammarlens}[title={What you've built}]
One more word for describing what you lived and felt.
\end{grammarlens}

\subsection*{Your turn}
\begin{itemize}
\raggedright
  \item \emph{Say it:} \emph{variṣṭha}
  \item \emph{Say it:} \emph{variṣṭha}, once more
  \item \emph{Say it:} \emph{variṣṭha adhikārī āj baiṭhkīlā āle.} — The senior officers came to the meeting today.
  \item \emph{Recall:} say the Marathi for a promotion, then the Marathi for a position, then say \emph{variṣṭha} again
\end{itemize}

\begin{tcolorbox}[breakable,colback=teal!4,colframe=teal!35!black,title={Before you move on}]
What does \mr{वरिष्ठ} mean? (\textquotedblleft{}Senior\textquotedblright{}.) Say it once more.
\end{tcolorbox}
\section[\texorpdfstring{\mr{कनिष्ठ} (kaniṣṭha)}{कनिष्ठ (kaniṣṭha)}]{\mr{कनिष्ठ} (kaniṣṭha) — junior}

\label{lesson:MR-C509-kanistha}



\begin{quote}
Before the new one: say the Marathi for a position, then the Marathi for senior.
\end{quote}

\subsection*{You'll want to know: \mr{कनिष्ठ}}
\textbf{\mr{कनिष्ठ}} — \emph{kaniṣṭha} — \textquotedblleft{}junior\textquotedblright{}.

\begin{quote}
\mr{कनिष्ठ} \mr{कर्मचारी} \mr{वरिष्ठांकडून} \mr{काम} \mr{शिकतात}.
\end{quote}

\begin{quote}
\emph{kaniṣṭha karmacārī variṣṭhāṅkaḍūn kām śiktāt.}
\end{quote}

\begin{quote}
Junior staff learn the work from their seniors.
\end{quote}

\begin{cousinweb}
\textbf{\mr{कनिष्ठ}} (\emph{kaniṣṭha}) is a tatsama, taken from Sanskrit \emph{kaniṣṭha}, \textquotedblleft{}youngest, least\textquotedblright{} (Monier-Williams).
\end{cousinweb}

\begin{grammarlens}[title={What you've built}]
One more word for describing what you lived and felt.
\end{grammarlens}

\subsection*{Your turn}
\begin{itemize}
\raggedright
  \item \emph{Say it:} \emph{kaniṣṭha}
  \item \emph{Say it:} \emph{kaniṣṭha}, once more
  \item \emph{Say it:} \emph{kaniṣṭha karmacārī variṣṭhāṅkaḍūn kām śiktāt.} — Junior staff learn the work from their seniors.
  \item \emph{Recall:} say the Marathi for a position, then the Marathi for senior, then say \emph{kaniṣṭha} again
\end{itemize}

\begin{tcolorbox}[breakable,colback=teal!4,colframe=teal!35!black,title={Before you move on}]
What does \mr{कनिष्ठ} mean? (\textquotedblleft{}Junior\textquotedblright{}.) Say it once more.
\end{tcolorbox}
\section[\texorpdfstring{\mr{झटणे} (jhaṭṇe)}{झटणे (jhaṭṇe)}]{\mr{झटणे} (jhaṭṇe) — to strive; to work hard}

\label{lesson:MR-C509-jhatne}



\begin{quote}
Before the new one: say the Marathi for senior, then the Marathi for junior.
\end{quote}

\subsection*{You'll want to know: \mr{झटणे}}
\textbf{\mr{झटणे}} — \emph{jhaṭṇe} — \textquotedblleft{}to strive; to work hard\textquotedblright{}.

\begin{quote}
\mr{पदोन्नतीसाठी} \mr{तो} \mr{दिवसरात्र} \mr{झटला}.
\end{quote}

\begin{quote}
\emph{padonnatīsāṭhī to divasrātra jhaṭlā.}
\end{quote}

\begin{quote}
He strove day and night for the promotion.
\end{quote}

Past (he …): \mr{तो} \mr{झटला}, \emph{to jhaṭlā}.

The goal takes -\mr{साठी}: \mr{यशासाठी} \mr{झटणे}.

\begin{cousinweb}
\textbf{\mr{झटणे}} (\emph{jhaṭṇe}) is an expressive Marathi verb. CDIAL lists a Marathi \emph{jhaṭṇe}, \textquotedblleft{}to collide violently\textquotedblright{}, under \emph{*jhaṭ}, \textquotedblleft{}sudden movement\textquotedblright{} (CDIAL 5327), but not the sense \textquotedblleft{}to strive\textquotedblright{}, so the link is not settled.
\end{cousinweb}

\begin{grammarlens}[title={What you've built}]
One more word for describing what you lived and felt.
\end{grammarlens}

\subsection*{Your turn}
\begin{itemize}
\raggedright
  \item \emph{Say it:} \emph{jhaṭṇe}
  \item \emph{Say it:} \emph{jhaṭṇe}, once more
  \item \emph{Say it:} \emph{padonnatīsāṭhī to divasrātra jhaṭlā.} — He strove day and night for the promotion.
  \item \emph{Recall:} say the Marathi for senior, then the Marathi for junior, then say \emph{jhaṭṇe} again
\end{itemize}

\begin{tcolorbox}[breakable,colback=teal!4,colframe=teal!35!black,title={Before you move on}]
What does \mr{झटणे} mean? (\textquotedblleft{}To strive; to work hard\textquotedblright{}.) Say it once more.
\end{tcolorbox}
